% !TEX root = ../Thesis.tex

This appendix lists the sources on which the thesis draws, grouped by the role
they play in the work rather than alphabetically, together with the chapter
that uses each one. Entries marked \emph{(preprint)} are cited from the version
consulted; where a later peer-reviewed version exists, the published reference
should replace it before submission.

\section{Causal inference and CATE estimation}

\begin{enumerate}
  \item Kennedy, E. H. \emph{Towards Optimal Doubly Robust Estimation of
    Heterogeneous Causal Effects.} arXiv:2004.14497v5 (preprint).
    \\\textit{Used in Ch.~2:} the doubly robust pseudo-outcome of
    Eq.~\eqref{eq:aipw-pseudo-outcome} and the two-stage DR-learner
    construction.

  \item Kennedy, E. H. \emph{Semiparametric Doubly Robust Targeted Double
    Machine Learning: A Review.} arXiv:2203.06469v2 [stat.ME], 26 January 2023
    (preprint).
    \\\textit{Used in Ch.~2:} influence-function framing of the AIPW estimator
    and the double-machine-learning nuisance structure behind the causal
    forest.

  \item Salditt, M., Eckes, T., Nestler, S. \emph{A Tutorial Introduction to
    Heterogeneous Treatment Effect Estimation with Meta-learners.}
    Administration and Policy in Mental Health and Mental Health Services
    Research, 51:650--673, 2024. DOI: 10.1007/s10488-023-01303-9.
    \\\textit{Used in Ch.~2:} the meta-learner taxonomy, in particular the
    T-learner and the S-learner families.
\end{enumerate}

\section{Validation of treatment prioritization rules}

The evaluation machinery of Section~2.3 --- the TOC curve, its AUTOC and QINI
summaries, and the group calibration test --- is applied through the
\texttt{econml} implementation (\texttt{econml.validate.DRTester}), which rests
on the following sources.

\begin{enumerate}\setcounter{enumi}{3}
  \item Chernozhukov, V. et al. \emph{Generic Machine Learning Inference on
    Heterogeneous Treatment Effects in Randomized Experiments.}
    arXiv:1712.04802, 2022 (preprint).
    \\\textit{Used in Ch.~2:} the group average treatment effect (GATES) and
    best-linear-predictor tests underlying group calibration.

  \item Dwivedi, R. et al. \emph{Stable Discovery of Interpretable Subgroups via
    Calibration in Causal Studies.} arXiv:2008.10109, 2020 (preprint).
    \\\textit{Used in Ch.~2:} the calibration $R^2$ normalized against the
    constant-effect baseline.

  \item Radcliffe, N. \emph{Using control groups to target on predicted lift:
    building and assessing uplift models.} Direct Marketing Analytics Journal,
    pages 14--21, 2007.
    \\\textit{Used in Ch.~2:} the QINI curve and its group-size weighting.
\end{enumerate}

\section{Adaptive enrolment, active and online learning}

\begin{enumerate}\setcounter{enumi}{6}
  \item Kato, M., Oga, A., Komatsubara, W., Inokuchi, R. \emph{Active Adaptive
    Experimental Design for Treatment Effect Estimation with Covariate Choice.}
    arXiv:2403.03589v2 (preprint).
    \\\textit{Used in Ch.~3 and Ch.~8:} joint optimization of covariate density
    and propensity score as the reference framework for guided enrolment.

  \item Cacciarelli, D., Kulahci, M. \emph{Active learning for data streams: a
    survey.} Machine Learning (Springer), 2023.
    DOI: 10.1007/s10994-023-06454-2.
    \\\textit{Used in Ch.~3:} the stream-based versus pool-based distinction
    that frames sequential patient enrolment as an online problem.

  \item Cacciarelli, D., Kulahci, M., Tyssedal, J. S. \emph{Stream-based active
    learning with linear models.} (preprint).
    \\\textit{Used in Ch.~3:} informativeness thresholding on a stream of
    incoming units.

  \item Cacciarelli, D., Kulahci, M., Tyssedal, J. \emph{Online Active Learning
    for Soft Sensor Development using Semi-Supervised Autoencoders.} ICML 2022
    Workshop on Adaptive Experimental Design and Active Learning in the Real
    World.
    \\\textit{Used in Ch.~3:} adaptation of pool-based query strategies to the
    stream-based setting.

  \item Fontaine, X., Perrault, P., Valko, M., Perchet, V. \emph{Online
    A-Optimal Design and Active Linear Regression.} arXiv:1906.08509v2
    [stat.ML], 30 December 2020 (preprint).
    \\\textit{Used in Ch.~3:} optimal-experiment-design view of budget
    allocation under heteroscedasticity.

  \item Riquelme, C., Johari, R., Zhang, B. \emph{Online Active Linear
    Regression via Thresholding.} arXiv:1602.02845v4 [stat.ML], 21 December
    2016 (preprint).
    \\\textit{Used in Ch.~3:} threshold-based selection of the most informative
    observations under a fixed experimentation budget.
\end{enumerate}

\section{Trial and dataset documentation}

\begin{enumerate}\setcounter{enumi}{12}
  \item Valgimigli, M., Frigoli, E., Heg, D., Tijssen, J., J\"uni, P., Vranckx,
    P., et al., for the MASTER DAPT Investigators. \emph{Dual Antiplatelet
    Therapy after PCI in Patients at High Bleeding Risk.} New England Journal
    of Medicine, 385(18):1643--1655, 28 October 2021.
    DOI: 10.1056/NEJMoa2108749.
    \\\textit{Used in Ch.~4 and Ch.~7:} trial design, randomization at the
    one-month landmark visit, endpoint definitions and headline results.

  \item \emph{Supplementary Appendix} to the above.
    \\\textit{Used in Ch.~4:} high-bleeding-risk criteria, inclusion and
    exclusion criteria, outcome definitions and the per-protocol population.

  \item Maroni, G., Azzimonti, L. \emph{Machine Learning Prediction of Clinical
    Outcomes in MASTER DAPT. Draft report: predictive performance, SHAP
    interpretation, and variable dictionary.} Internal report, 6 May 2026.
    Unpublished.
    \\\textit{Used in Ch.~5:} baseline predictive performance, SHAP-based
    interpretation and the variable dictionary of the modelling dataset.
\end{enumerate}
